% Blueprint: Brouwer fixed-point theorem
% Created by Numina
%
% Strategy.  Mathlib has neither Brouwer nor the usual high-level tools used to
% prove it (it lacks singular homology of spheres, excision / Mayer--Vietoris
% for spaces, the no-retraction theorem, and topological degree).  We therefore
% take the elementary combinatorial route through Sperner's lemma, every leaf of
% which bottoms out in Mathlib's combinatorics, convex geometry, and point-set
% topology:
%
%   Sperner's lemma  -->  Brouwer for a full-dimensional simplex
%                    -->  Brouwer for the closed ball (Mathlib convex homeo)
%                    -->  Brouwer for an arbitrary compact convex set.
%
% Throughout, E_n := EuclideanSpace R (Fin n) and
% S_n := convexHull {0, e_1, ..., e_n} \subseteq E_n is the standard
% full-dimensional corner simplex.

\begin{document}

\section{The fixed-point property}

\begin{definition}[Fixed-point property]
    \label{def:fpp}
    \lean{LeanEval.Topology.HasFPP}
    \leanfile{LeanEval/Topology/Brouwer/FixedPoint.lean}
    \leanok
    A topological space $X$ has the \emph{fixed-point property} (FPP) if every
    continuous map $f : X \to X$ has a fixed point: there exists $x \in X$ with
    $f(x) = x$.
\end{definition}

\begin{lemma}[FPP is a topological invariant]
    \label{lem:fpp-homeo}
    \lean{LeanEval.Topology.fpp_homeo}
    \leanfile{LeanEval/Topology/Brouwer/FixedPoint.lean}
    \uses{def:fpp}
    \leanok
    If $X$ and $Y$ are homeomorphic and $X$ has the fixed-point property, then
    $Y$ has the fixed-point property.
\end{lemma}
\begin{proof}
    \uses{def:fpp}
    Let $h : X \simeq Y$ be a homeomorphism and let $g : Y \to Y$ be continuous.
    Then $h^{-1} \circ g \circ h : X \to X$ is continuous, so it has a fixed
    point $x$. Setting $y = h(x)$ and applying $h$ to $h^{-1}(g(h(x))) = x$ gives
    $g(y) = y$.
\end{proof}

\begin{lemma}[Subtype FPP gives an ambient fixed point]
    \label{lem:fpp-ambient}
    \lean{LeanEval.Topology.fpp_ambient}
    \leanfile{LeanEval/Topology/Brouwer/FixedPoint.lean}
    \uses{def:fpp}
    \leanok
    Let $K \subseteq E_d$ and suppose the subspace $K$ (with the subspace
    topology) has the fixed-point property. Then for every $f : E_d \to E_d$
    that is continuous on $K$ and maps $K$ into $K$, there exists $x \in K$ with
    $f(x) = x$.
\end{lemma}
\begin{proof}
    \uses{def:fpp}
    The hypotheses package $f$ as a continuous self-map $\tilde f : K \to K$ of
    the subspace. By the fixed-point property of $K$ there is $x \in K$ with
    $\tilde f(x) = x$, i.e.\ $f(x) = x$.
\end{proof}

\section{Abstract Sperner combinatorics}

% The finite combinatorics of Sperner's lemma is developed here over an abstract
% axiomatized incidence model, with no geometry.  All planar/PL content is
% supplied later by the geometric layer, which merely has to exhibit a model and
% verify the balancedness axiom.

\begin{definition}[Abstract labeled incidence model]
    \label{def:incidence-model}
    \lean{LeanEval.Topology.IncidenceModel}
    \leanfile{LeanEval/Topology/Brouwer/Incidence.lean}
    \leanok
    A \emph{labeled incidence model} of dimension $n$ consists of a finite type
    of \emph{cells} and a finite type of \emph{facets}, an \emph{incidence}
    relation $\mathrm{Inc}$ between cells and facets, a \emph{boundary} predicate
    $\partial$ on facets, and a labeling assigning to each cell a multiset of
    $n+1$ labels from $\{0, \dots, n\}$ and to each facet a multiset of $n$
    labels from $\{0, \dots, n\}$, subject to the compatibility axiom that
    whenever $\mathrm{Inc}(C, F)$ the facet multiset of $F$ is obtained from the
    cell multiset of $C$ by deleting one label. A facet is a \emph{door} if its
    label multiset is exactly $\{0, \dots, n-1\}$ (each label once); a cell is
    \emph{rainbow} if its label multiset is exactly $\{0, \dots, n\}$ (each label
    once). A door is an \emph{interior door} if $\neg\partial$ holds of it and a
    \emph{boundary door} if $\partial$ holds.
\end{definition}

\begin{definition}[Balanced model]
    \label{def:balanced-model}
    \lean{LeanEval.Topology.IncidenceModel.Balanced}
    \leanfile{LeanEval/Topology/Brouwer/Incidence.lean}
    \uses{def:incidence-model}
    \leanok
    A labeled incidence model is \emph{balanced} if every facet with
    $\neg\partial$ (an interior facet) is incident to exactly two cells and every
    facet with $\partial$ (a boundary facet) is incident to exactly one cell.
\end{definition}

\begin{lemma}[Rainbow cells have exactly one door]
    \label{lem:cell-door-rainbow}
    \lean{LeanEval.Topology.cell_door_rainbow}
    \leanfile{LeanEval/Topology/Brouwer/Incidence.lean}
    \uses{def:incidence-model}
    A rainbow cell is incident to exactly one door.
\end{lemma}
\begin{proof}
    \uses{def:incidence-model}
    The cell's label multiset is $\{0, \dots, n\}$, each label appearing once.
    A facet incident to it is obtained by deleting one label, and the result is
    $\{0, \dots, n-1\}$ exactly when the deleted label is $n$. The vertex labeled
    $n$ is unique, so exactly one incident facet is a door
    (\texttt{Finset.card\_eq\_one}).
\end{proof}

\begin{lemma}[Non-rainbow cells have an even number of doors]
    \label{lem:cell-door-nonrainbow}
    \lean{LeanEval.Topology.cell_door_nonrainbow}
    \leanfile{LeanEval/Topology/Brouwer/Incidence.lean}
    \uses{def:incidence-model}
    A cell that is not rainbow is incident to either $0$ or $2$ doors; in
    particular an even number.
\end{lemma}
\begin{proof}
    \uses{def:incidence-model}
    Deleting one label of the cell yields $\{0, \dots, n-1\}$ only if the cell's
    underlying label set already contains $\{0, \dots, n-1\}$. As the cell is not
    rainbow, its $n+1$ labels are not all distinct, so the only remaining
    possibility is that its label set is exactly $\{0, \dots, n-1\}$ with one
    label repeated; then deleting either copy of the repeated label gives a door,
    and no other deletion does, so there are exactly $2$ doors. If instead some
    label of $\{0, \dots, n-1\}$ is missing from the cell, no deletion yields a
    door and there are $0$.
\end{proof}

\begin{lemma}[Per-cell door parity]
    \label{lem:cell-door-count}
    \lean{LeanEval.Topology.cell_door_count}
    \leanfile{LeanEval/Topology/Brouwer/Incidence.lean}
    \uses{lem:cell-door-rainbow, lem:cell-door-nonrainbow}
    \leanok
    The number of doors incident to a cell is $1$ if the cell is rainbow and is
    even otherwise; in particular it is odd if and only if the cell is rainbow.
\end{lemma}
\begin{proof}
    \uses{lem:cell-door-rainbow, lem:cell-door-nonrainbow}
    Immediate from \Cref{lem:cell-door-rainbow} (rainbow case: exactly $1$, which
    is odd) and \Cref{lem:cell-door-nonrainbow} (non-rainbow case: even).
\end{proof}

\begin{lemma}[Door count has rainbow parity]
    \label{lem:door-count-rainbow-parity}
    \lean{LeanEval.Topology.door_count_rainbow_parity}
    \leanfile{LeanEval/Topology/Brouwer/Incidence.lean}
    \uses{lem:cell-door-count}
    \leanok
    In any labeled incidence model, $\sum_{C} \#\{\text{doors incident to } C\}$
    has the same parity as the number of rainbow cells.
\end{lemma}
\begin{proof}
    \uses{lem:cell-door-count}
    By \Cref{lem:cell-door-count} each summand is odd exactly for the rainbow
    cells, so modulo $2$ the sum equals the number of cells with an odd summand,
    i.e.\ the number of rainbow cells (sum of a $\{0,1\}$-valued indicator,
    \texttt{Finset.card\_filter} and \texttt{Nat.even\_add}).
\end{proof}

\begin{lemma}[Door count via incidence]
    \label{lem:door-count-incidence}
    \lean{LeanEval.Topology.door_count_incidence}
    \leanfile{LeanEval/Topology/Brouwer/Incidence.lean}
    \uses{def:balanced-model}
    \leanok
    In a balanced model,
    $\sum_{C} \#\{\text{doors incident to } C\}
       = 2\,\#\{\text{interior doors}\} + \#\{\text{boundary doors}\}$;
    in particular it has the same parity as the number of boundary doors.
\end{lemma}
\begin{proof}
    \uses{def:balanced-model}
    Double count incident (cell, door) pairs: by the bipartite double-counting
    identity \texttt{Finset.sum\_card\_bipartiteAbove\_eq\_sum\_card\_bipartiteBelow},
    $\sum_{C} \#\{\text{doors } F : \mathrm{Inc}(C,F)\}
       = \sum_{\text{doors } F} \#\{C : \mathrm{Inc}(C,F)\}$.
    By \Cref{def:balanced-model} the inner count is $2$ for an interior door and
    $1$ for a boundary door, giving
    $2\,\#\{\text{interior doors}\} + \#\{\text{boundary doors}\}$, which is
    congruent to $\#\{\text{boundary doors}\}$ modulo $2$.
\end{proof}

\begin{lemma}[Boundary parity count]
    \label{lem:sperner-parity}
    \lean{LeanEval.Topology.sperner_parity}
    \leanfile{LeanEval/Topology/Brouwer/Incidence.lean}
    \uses{lem:door-count-rainbow-parity, lem:door-count-incidence}
    \leanok
    In a balanced labeled incidence model, the number of rainbow cells has the
    same parity as the number of boundary doors.
\end{lemma}
\begin{proof}
    \uses{lem:door-count-rainbow-parity, lem:door-count-incidence}
    The single quantity $\sum_{C} \#\{\text{doors incident to } C\}$ is congruent
    modulo $2$ both to the number of rainbow cells
    (\Cref{lem:door-count-rainbow-parity}) and to the number of boundary doors
    (\Cref{lem:door-count-incidence}); hence those two numbers are congruent to
    each other.
\end{proof}

\section{Triangulations and their incidence model}

\begin{definition}[Triangulation]
    \label{def:triangulation}
    \lean{LeanEval.Topology.Triangulation}
    \leanfile{LeanEval/Topology/Brouwer/Triangulation.lean}
    \leanok
    A \emph{triangulation} of the corner simplex
    $S_n = \mathrm{convexHull}\,\{0, e_1, \dots, e_n\}$ is a bespoke structure
    -- richer than Mathlib's \texttt{Geometry.SimplicialComplex}, which records
    faces and the face-to-face condition but neither purity, dimension, nor the
    requirement that the underlying space be a prescribed set -- packaging the
    following data and axioms:
    \begin{itemize}
        \item a finite set of \emph{cells}, each cell a \texttt{Finset} of points
            of $E_n$ that is affinely independent (so its convex hull is a
            nondegenerate simplex), of dimension at most $n$;
        \item \emph{closure under faces}: every nonempty subset of a cell is a
            cell;
        \item the \emph{face-to-face} condition: for cells $C, C'$ the
            intersection $\mathrm{conv}(C) \cap \mathrm{conv}(C')$ is the convex
            hull of $C \cap C'$ (a common face);
        \item \emph{purity of dimension $n$}: every cell is a subset of a
            \emph{maximal cell} with exactly $n+1$ vertices;
        \item the \emph{covering} axiom: the union of the convex hulls of the
            maximal cells equals $S_n$.
    \end{itemize}
    The \emph{maximal cells} are the cells with $n+1$ vertices and the
    \emph{$(n-1)$-faces} (or \emph{facets}) are the cells with $n$ vertices. This
    structure cannot be expressed by Mathlib's single-simplex constructors,
    because the vertices of distinct cells are affinely independent only
    cell-by-cell, not jointly; the bespoke API above is the shared foundational
    layer used throughout this section.
\end{definition}

\begin{definition}[Carrier face]
    \label{def:carrier-face}
    \lean{LeanEval.Topology.carrierFace}
    \leanfile{LeanEval/Topology/Brouwer/Triangulation.lean}
    \uses{def:triangulation}
    \leanok
    Writing a point $x \in S_n$ in barycentric coordinates $(x_0, \dots, x_n)$
    relative to the corner vertices $0, e_1, \dots, e_n$ (so $x_i \ge 0$ and
    $\sum_i x_i = 1$), the \emph{carrier face} of $x$ is the closed face of $S_n$
    spanned by the corner vertices indexed by
    $\mathrm{supp}(x) = \{\, i : x_i > 0 \,\}$; it is the smallest face of $S_n$
    containing $x$.
\end{definition}

\begin{definition}[Labeling and Sperner labeling]
    \label{def:sperner-labeling}
    \lean{LeanEval.Topology.IsSpernerLabeling}
    \leanfile{LeanEval/Topology/Brouwer/Triangulation.lean}
    \uses{def:triangulation, def:carrier-face}
    \leanok
    A \emph{labeling} of a triangulation $T$ of $S_n$ is a function $\ell$ from
    its vertex set to $\{0, 1, \dots, n\}$. It is a \emph{Sperner labeling} if
    $\ell(v) \in \mathrm{supp}(v)$ for every vertex $v$, i.e.\ each vertex
    receives a label belonging to the index set of its carrier face. An
    $(n-1)$-face of $T$ is a \emph{door} if its $n$ vertices carry exactly the
    labels $\{0, \dots, n-1\}$. A maximal cell is \emph{rainbow} if its $n+1$
    vertices carry all $n+1$ distinct labels $\{0, \dots, n\}$.
\end{definition}

\begin{definition}[Incidence model of a triangulation]
    \label{def:triangulation-model}
    \lean{LeanEval.Topology.triangulationModel}
    \leanfile{LeanEval/Topology/Brouwer/Triangulation.lean}
    \uses{def:triangulation, def:sperner-labeling, def:incidence-model}
    \leanok
    Given a Sperner-labeled triangulation $T$ of $S_n$, its associated
    \emph{labeled incidence model} (\Cref{def:incidence-model}) has the maximal
    cells of $T$ as cells, the $(n-1)$-faces of $T$ as facets, incidence given by
    ``is a face of'' (i.e.\ $\mathrm{Inc}(C, F)$ iff $F \subseteq C$), boundary
    predicate $\partial F$ given by $\mathrm{conv}(F) \subseteq \partial S_n$, and
    labels the vertex labels of $T$. The facet multiset of $F$ is the multiset of
    labels of its $n$ vertices, which depends only on $F$; deleting from a maximal
    cell $C \supseteq F$ the label of the unique vertex of $C$ outside $F$
    recovers it, so the compatibility axiom holds. Doors and rainbow cells of the
    model coincide with the geometric doors and rainbow maximal cells of $T$.
\end{definition}

\begin{lemma}[A facet spans an affine hyperplane]
    \label{lem:facet-hyperplane}
    \lean{LeanEval.Topology.facet_hyperplane}
    \leanfile{LeanEval/Topology/Brouwer/Triangulation.lean}
    \uses{def:triangulation}
    The affine span $H$ of the $n$ vertices of an $(n-1)$-face $F$ of a
    triangulation of $S_n$ is an affine hyperplane of $E_n$ (codimension $1$),
    and $E_n \setminus H$ has exactly two connected components, the two open
    half-spaces $H^+, H^-$ bounded by $H$.
\end{lemma}
\begin{proof}
    \uses{def:triangulation}
    The $n$ vertices of $F$ are affinely independent, so $H$ is an affine
    subspace whose direction has dimension $n-1$, a hyperplane. The complement of
    an affine hyperplane in $E_n$ is the disjoint union of the two open
    half-spaces cut out by any defining affine functional (\texttt{AffineMap}
    to $\mathbb R$ with the two strict sign conditions).
\end{proof}

\begin{lemma}[An incident cell lies on one side of its facet]
    \label{lem:cell-one-side}
    \lean{LeanEval.Topology.cell_one_side}
    \leanfile{LeanEval/Topology/Brouwer/Triangulation.lean}
    \uses{lem:facet-hyperplane, def:triangulation}
    If a maximal cell $C$ has the $(n-1)$-face $F$ as a face, its remaining
    vertex $w$ lies off the hyperplane $H = \mathrm{aff}(F)$, and $C$ is contained
    in the closed half-space on $w$'s side, with the relative interior of $C$ in
    the corresponding open half-space.
\end{lemma}
\begin{proof}
    \uses{lem:facet-hyperplane, def:triangulation}
    Write $C = F \cup \{w\}$. Affine independence of the $n+1$ vertices of $C$
    forces $w \notin H$, so $w \in H^+$ say. As $\mathrm{conv}(C)$ is the convex
    hull of points of $H \cup \{w\}$, all in the closed half-space $\overline{H^+}$,
    convexity gives $\mathrm{conv}(C) \subseteq \overline{H^+}$, and any point with
    positive $w$-barycentric weight lies in $H^+$.
\end{proof}

\begin{lemma}[At most one incident cell per side]
    \label{lem:one-cell-per-side}
    \lean{LeanEval.Topology.one_cell_per_side}
    \leanfile{LeanEval/Topology/Brouwer/Triangulation.lean}
    \uses{lem:cell-one-side, def:triangulation}
    Two distinct maximal cells both having $F$ as a face lie on opposite sides of
    $H = \mathrm{aff}(F)$; equivalently each open side $H^\pm$ contains the
    relative interior of at most one maximal cell incident to $F$.
\end{lemma}
\begin{proof}
    \uses{lem:cell-one-side, def:triangulation}
    Suppose maximal cells $C \ne C'$ both contain $F$ and, by
    \Cref{lem:cell-one-side}, both have relative interior in $H^+$. By the
    face-to-face axiom $\mathrm{conv}(C) \cap \mathrm{conv}(C') =
    \mathrm{conv}(C \cap C')$ is a common face containing $F$, hence equal to $F$
    or to a cell of dimension $\ge n$. The latter forces $C = C'$; the former is
    impossible because two distinct $n$-simplices sharing the facet $F$ and lying
    on the same side overlap in a neighborhood of a relative-interior point of
    $F$ beyond $F$ itself. So at most one incident maximal cell occupies each
    open side.
\end{proof}

\begin{lemma}[Both sides of an interior facet are occupied]
    \label{lem:interior-facet-both-sides}
    \lean{LeanEval.Topology.interior_facet_both_sides}
    \leanfile{LeanEval/Topology/Brouwer/Triangulation.lean}
    \uses{lem:facet-hyperplane, def:triangulation}
    If an $(n-1)$-face $F$ is not contained in $\partial S_n$, then a
    relative-interior point of $F$ lies in the topological interior of $S_n$, and
    each open side $H^\pm$ contains the relative interior of some maximal cell
    incident to $F$.
\end{lemma}
\begin{proof}
    \uses{lem:facet-hyperplane, def:triangulation}
    Pick $p$ in the relative interior of $F$. Since $F \not\subseteq \partial
    S_n$, $p$ is an interior point of $S_n$, so a small ball $B$ around $p$ lies
    in $S_n$ and meets both open half-spaces $H^\pm$. By the covering axiom each
    point of $B$ on a given side lies in some maximal cell; a point close enough
    to $p$ and strictly on one side lies in a maximal cell having $F$ as a face on
    that side. Thus both sides are occupied.
\end{proof}

\begin{lemma}[A boundary facet is occupied on one side only]
    \label{lem:boundary-facet-one-side}
    \lean{LeanEval.Topology.boundary_facet_one_side}
    \leanfile{LeanEval/Topology/Brouwer/Triangulation.lean}
    \uses{lem:facet-hyperplane, def:triangulation}
    If an $(n-1)$-face $F$ satisfies $\mathrm{conv}(F) \subseteq \partial S_n$,
    then near a relative-interior point of $F$ the set $S_n$ lies in one closed
    half-space of $H = \mathrm{aff}(F)$, so only one open side carries the
    relative interior of a maximal cell incident to $F$.
\end{lemma}
\begin{proof}
    \uses{lem:facet-hyperplane, def:triangulation}
    A facet of $S_n$ lies in a supporting affine hyperplane of the convex set
    $S_n$, which (where it meets the relative interior of $F$) coincides with
    $H$; on one side of a supporting hyperplane $S_n$ is empty. Hence no maximal
    cell has relative interior on that side, and by the covering axiom the other
    side is occupied.
\end{proof}

\begin{lemma}[Door incidence]
    \label{lem:door-incidence}
    \lean{LeanEval.Topology.door_incidence}
    \leanfile{LeanEval/Topology/Brouwer/Triangulation.lean}
    \uses{lem:interior-facet-both-sides, lem:one-cell-per-side, lem:boundary-facet-one-side, lem:cell-one-side}
    In a triangulation of $S_n$, an $(n-1)$-face not contained in the topological
    boundary $\partial S_n$ is a face of exactly two maximal cells, while an
    $(n-1)$-face contained in $\partial S_n$ is a face of exactly one maximal
    cell.
\end{lemma}
\begin{proof}
    \uses{lem:interior-facet-both-sides, lem:one-cell-per-side, lem:boundary-facet-one-side, lem:cell-one-side}
    Every incident maximal cell occupies exactly one open side of $H$
    (\Cref{lem:cell-one-side}), and each side carries at most one
    (\Cref{lem:one-cell-per-side}). For an interior facet both sides are occupied
    (\Cref{lem:interior-facet-both-sides}), giving exactly two; for a boundary
    facet only one side is (\Cref{lem:boundary-facet-one-side}), giving exactly
    one.
\end{proof}

\begin{lemma}[The triangulation model is balanced]
    \label{lem:triangulation-model-balanced}
    \lean{LeanEval.Topology.triangulation_model_balanced}
    \leanfile{LeanEval/Topology/Brouwer/Triangulation.lean}
    \uses{def:triangulation-model, lem:door-incidence, def:balanced-model}
    \leanok
    The labeled incidence model of a Sperner-labeled triangulation of $S_n$ is
    balanced.
\end{lemma}
\begin{proof}
    \uses{def:triangulation-model, lem:door-incidence}
    The boundary predicate of the model is exactly $\mathrm{conv}(F) \subseteq
    \partial S_n$, so balancedness -- interior facets incident to two cells and
    boundary facets to one -- is precisely the statement of
    \Cref{lem:door-incidence}.
\end{proof}

\begin{lemma}[The bottom facet is a triangulation of $S_{n-1}$]
    \label{lem:facet-is-triangulation}
    \lean{LeanEval.Topology.facet_is_triangulation}
    \leanfile{LeanEval/Topology/Brouwer/Triangulation.lean}
    \uses{def:triangulation}
    Let $T$ be a triangulation of $S_n$ with $n \ge 1$. The cells of $T$
    contained in $\{x_n = 0\}$, transported by the affine isomorphism dropping the
    last coordinate, form a triangulation of $S_{n-1}$.
\end{lemma}
\begin{proof}
    \uses{def:triangulation}
    The set $\{x_n = 0\} \cap S_n$ is the closed face spanned by
    $0, e_1, \dots, e_{n-1}$, affinely isomorphic to $S_{n-1}$. Restricting a
    triangulation to a closed face of $S_n$ preserves affine independence,
    closure under faces, and the face-to-face condition; purity and covering for
    the face follow because every point of the facet lies in a maximal cell of
    $T$ whose intersection with the facet is a maximal $(n-1)$-cell of the face.
    Transporting by the affine isomorphism gives a triangulation of $S_{n-1}$.
\end{proof}

\begin{lemma}[The restricted labeling is Sperner]
    \label{lem:facet-labeling-sperner}
    \lean{LeanEval.Topology.facet_labeling_sperner}
    \leanfile{LeanEval/Topology/Brouwer/Triangulation.lean}
    \uses{lem:facet-is-triangulation, def:sperner-labeling, def:carrier-face}
    The labels of $T$ restrict, on the triangulation of $S_{n-1}$ of
    \Cref{lem:facet-is-triangulation}, to a Sperner labeling, with all labels
    lying in $\{0, \dots, n-1\}$.
\end{lemma}
\begin{proof}
    \uses{lem:facet-is-triangulation, def:sperner-labeling, def:carrier-face}
    A vertex $v$ in $\{x_n = 0\}$ has $x_n(v) = 0$, so $n \notin \mathrm{supp}(v)$,
    and the Sperner condition $\ell(v) \in \mathrm{supp}(v)$ forces
    $\ell(v) \in \{0, \dots, n-1\}$. The carrier-face condition relative to
    $S_{n-1}$ is the restriction of that for $S_n$, so the labeling is Sperner.
\end{proof}

\begin{lemma}[Restriction to a boundary facet]
    \label{lem:facet-restriction}
    \lean{LeanEval.Topology.facet_restriction}
    \leanfile{LeanEval/Topology/Brouwer/Triangulation.lean}
    \uses{lem:facet-is-triangulation, lem:facet-labeling-sperner}
    \leanok
    Let $T$ be a Sperner-labeled triangulation of $S_n$ with $n \ge 1$. The faces
    of $T$ contained in $\{x_n = 0\}$, identified with $S_{n-1}$, form a
    triangulation of $S_{n-1}$ whose inherited labeling is Sperner.
\end{lemma}
\begin{proof}
    \uses{lem:facet-is-triangulation, lem:facet-labeling-sperner}
    Combine \Cref{lem:facet-is-triangulation} (the restricted complex is a
    triangulation of $S_{n-1}$) with \Cref{lem:facet-labeling-sperner} (its
    inherited labeling is Sperner).
\end{proof}

\begin{lemma}[Boundary doors lie on the bottom facet]
    \label{lem:boundary-doors-bottom}
    \lean{LeanEval.Topology.boundary_doors_bottom}
    \leanfile{LeanEval/Topology/Brouwer/Triangulation.lean}
    \uses{def:sperner-labeling}
    Every door of $T$ contained in $\partial S_n$ lies on the facet
    $\{x_n = 0\}$.
\end{lemma}
\begin{proof}
    \uses{def:sperner-labeling}
    The boundary facets of $S_n$ are the faces $\{x_j = 0\}$ opposite the corner
    vertices $j \in \{0, \dots, n\}$. On $\{x_j = 0\}$ every vertex $v$ has
    $x_j(v) = 0$, so $j \notin \mathrm{supp}(v)$ and the Sperner condition gives
    $\ell(v) \ne j$; thus label $j$ is absent from any cell on that facet. A door
    carries every label in $\{0, \dots, n-1\}$, so if it lies on $\{x_j = 0\}$
    then $j \notin \{0, \dots, n-1\}$, i.e.\ $j = n$. Hence the door lies on
    $\{x_n = 0\}$.
\end{proof}

\begin{lemma}[Bottom-facet doors are rainbow cells]
    \label{lem:bottom-facet-doors-rainbow}
    \lean{LeanEval.Topology.bottom_facet_doors_rainbow}
    \leanfile{LeanEval/Topology/Brouwer/Triangulation.lean}
    \uses{lem:facet-restriction, def:sperner-labeling}
    Under the identification of \Cref{lem:facet-restriction}, the doors of $T$
    lying on $\{x_n = 0\}$ are exactly the rainbow maximal cells of the induced
    Sperner-labeled triangulation of $S_{n-1}$.
\end{lemma}
\begin{proof}
    \uses{lem:facet-restriction, def:sperner-labeling}
    A door on $\{x_n = 0\}$ is an $(n-1)$-cell whose $n$ vertices carry exactly
    the labels $\{0, \dots, n-1\}$; under the identification this is precisely an
    $(n-1)$-dimensional maximal cell of the induced triangulation of $S_{n-1}$ all
    of whose $n$ vertices carry distinct labels $\{0, \dots, n-1\}$, i.e.\ a
    rainbow cell.
\end{proof}

\begin{lemma}[Boundary doors are the rainbow cells of the facet]
    \label{lem:boundary-doors-facet}
    \lean{LeanEval.Topology.boundary_doors_facet}
    \leanfile{LeanEval/Topology/Brouwer/Triangulation.lean}
    \uses{lem:boundary-doors-bottom, lem:bottom-facet-doors-rainbow}
    The doors of $T$ contained in $\partial S_n$ are exactly the rainbow maximal
    cells of the induced Sperner-labeled triangulation of $S_{n-1}$ on
    $\{x_n = 0\}$.
\end{lemma}
\begin{proof}
    \uses{lem:boundary-doors-bottom, lem:bottom-facet-doors-rainbow}
    By \Cref{lem:boundary-doors-bottom} every boundary door lies on $\{x_n = 0\}$,
    and by \Cref{lem:bottom-facet-doors-rainbow} the doors there are exactly the
    rainbow cells of the induced triangulation of $S_{n-1}$.
\end{proof}

\begin{lemma}[Sperner's lemma]
    \label{lem:sperner}
    \lean{LeanEval.Topology.sperner}
    \leanfile{LeanEval/Topology/Brouwer/Triangulation.lean}
    \uses{lem:sperner-parity, lem:triangulation-model-balanced, lem:boundary-doors-facet}
    Every Sperner-labeled triangulation of $S_n$ has an odd number of rainbow
    maximal cells; in particular it has at least one rainbow cell.
\end{lemma}
\begin{proof}
    \uses{lem:sperner-parity, lem:triangulation-model-balanced, lem:boundary-doors-facet}
    Induct on $n$. For $n = 0$ the single vertex is forced to carry label $0$,
    giving exactly one rainbow cell. For $n \ge 1$, apply \Cref{lem:sperner-parity}
    to the model of $T$, which is balanced by
    \Cref{lem:triangulation-model-balanced}: the number of rainbow cells has the
    same parity as the number of boundary doors. By \Cref{lem:boundary-doors-facet}
    that count equals the number of rainbow cells of the induced Sperner-labeled
    triangulation of $S_{n-1}$, which is odd by the induction hypothesis. Hence
    $S_n$ has an odd, so nonzero, number of rainbow cells.
\end{proof}

\section{Triangulations of small mesh}

\begin{definition}[Mesh]
    \label{def:mesh}
    \lean{LeanEval.Topology.Triangulation.mesh}
    \leanfile{LeanEval/Topology/Brouwer/Triangulation.lean}
    \uses{def:triangulation}
    \leanok
    The \emph{mesh} of a triangulation is the maximum diameter of its maximal
    cells.
\end{definition}

\begin{definition}[Barycentric subdivision]
    \label{def:barycentric-subdivision}
    \lean{LeanEval.Topology.barycentricSubdivision}
    \leanfile{LeanEval/Topology/Brouwer/Subdivision.lean}
    \uses{def:triangulation}
    The \emph{barycentric subdivision} of a triangulation $T$ of $S_n$ is the
    complex whose maximal cells are, for each maximal cell
    $\mathrm{conv}\{v_0, \dots, v_n\}$ of $T$ and each chain
    $F_0 \subsetneq \dots \subsetneq F_n$ of its faces ordered by inclusion, the
    convex hull of the barycenters $b(F_0), \dots, b(F_n)$ of that chain.
\end{definition}

\begin{lemma}[Chains of faces cover a simplex]
    \label{lem:bary-chain-cover}
    \lean{LeanEval.Topology.bary_chain_cover}
    \leanfile{LeanEval/Topology/Brouwer/Subdivision.lean}
    \uses{def:barycentric-subdivision}
    For a single maximal cell, the chain simplices
    $\mathrm{conv}\{b(F_0), \dots, b(F_n)\}$ over the strictly increasing flags
    $F_0 \subsetneq \dots \subsetneq F_n$ of its faces cover the cell and have
    pairwise disjoint relative interiors; there are $n!$ of them.
\end{lemma}
\begin{proof}
    \uses{def:barycentric-subdivision}
    A point of the cell, written in barycentric coordinates, determines (after
    sorting its coordinates) a unique maximal flag of faces and expresses the
    point as a convex combination of the corresponding barycenters; this is the
    classical order-complex description, giving the $n!$ flags indexed by
    permutations, a cover, and disjoint relative interiors.
\end{proof}

\begin{lemma}[Chain barycenters are affinely independent]
    \label{lem:bary-chain-indep}
    \lean{LeanEval.Topology.bary_chain_indep}
    \leanfile{LeanEval/Topology/Brouwer/Subdivision.lean}
    \uses{def:barycentric-subdivision}
    The barycenters $b(F_0), \dots, b(F_n)$ of a strictly increasing flag
    $F_0 \subsetneq \dots \subsetneq F_n$ are affinely independent, so their
    convex hull is a nondegenerate $n$-simplex.
\end{lemma}
\begin{proof}
    \uses{def:barycentric-subdivision}
    Each $b(F_k)$ involves exactly the vertices of $F_k$ with positive weight,
    and the supports $F_0 \subsetneq \dots \subsetneq F_n$ are strictly nested, so
    no $b(F_k)$ is an affine combination of the others; hence the family is
    affinely independent.
\end{proof}

\begin{lemma}[Subdivisions match face-to-face]
    \label{lem:bary-face-to-face}
    \lean{LeanEval.Topology.bary_face_to_face}
    \leanfile{LeanEval/Topology/Brouwer/Subdivision.lean}
    \uses{def:barycentric-subdivision}
    \leanok
    On a face shared by two maximal cells, the chain simplices induced from
    either side coincide, so the subdivisions of adjacent maximal cells agree on
    their common face.
\end{lemma}
\begin{proof}
    \uses{def:barycentric-subdivision}
    The barycenter of a face depends only on that face, not on the maximal cell
    containing it, and the flags lying inside the shared face are the same when
    read from either side; hence the chain simplices supported on the shared face
    are identical.
\end{proof}

\begin{lemma}[Barycentric subdivision covers the same space]
    \label{lem:barycentric-same-space}
    \lean{LeanEval.Topology.barycentric_same_space}
    \leanfile{LeanEval/Topology/Brouwer/Subdivision.lean}
    \uses{lem:bary-chain-cover, lem:bary-chain-indep, lem:bary-face-to-face, def:triangulation}
    \leanok
    The barycentric subdivision of a triangulation of $S_n$ is again a
    triangulation of $S_n$.
\end{lemma}
\begin{proof}
    \uses{lem:bary-chain-cover, lem:bary-chain-indep, lem:bary-face-to-face}
    The chain simplices are nondegenerate $n$-simplices
    (\Cref{lem:bary-chain-indep}), they cover each maximal cell with disjoint
    relative interiors (\Cref{lem:bary-chain-cover}), and the pieces of adjacent
    cells agree on shared faces (\Cref{lem:bary-face-to-face}); hence they satisfy
    the triangulation axioms with the same underlying space $S_n$.
\end{proof}

\begin{lemma}[Barycenter distance inequality]
    \label{lem:barycenter-dist-bound}
    \lean{LeanEval.Topology.barycenter_dist_bound}
    \leanfile{LeanEval/Topology/Brouwer/Subdivision.lean}
    Let $G$ be a face of dimension $d$ of a simplex and $F$ a nonempty face of
    $G$. Then $\lVert b(F) - b(G) \rVert \le \frac{d}{d+1}\,\mathrm{diam}(G)$,
    where $b(\cdot)$ denotes the barycenter (centroid) of the vertex set.
\end{lemma}
\begin{proof}
    The barycenter $b(G)$ is the centroid of the $d+1$ vertices of $G$
    (\texttt{Finset.centroid\_eq\_centerMass}), and $b(F)$ is a convex
    combination of vertices of $G$. By the convex-combination distance estimate
    \texttt{Convexity.dist\_sConvexComb\_right\_le}, $\lVert b(F) - b(G) \rVert$
    is at most $\max_v \lVert v - b(G) \rVert$ over vertices $v$ of $G$. For each
    vertex $v$, $v - b(G) = \frac{1}{d+1}\sum_{w}(v - w)$ averages the $d$ nonzero
    differences $v - w$, each of norm $\le \mathrm{diam}(G)$, giving
    $\lVert v - b(G)\rVert \le \frac{d}{d+1}\,\mathrm{diam}(G)$.
\end{proof}

\begin{lemma}[Barycentric subdivision shrinks the mesh]
    \label{lem:barycentric-shrink}
    \lean{LeanEval.Topology.barycentric_shrink}
    \leanfile{LeanEval/Topology/Brouwer/Subdivision.lean}
    \uses{def:mesh, def:barycentric-subdivision, lem:barycentric-same-space, lem:barycenter-dist-bound}
    The mesh of the barycentric subdivision of a triangulation of $S_n$ is at
    most $\frac{n}{n+1}$ times the mesh of the original.
\end{lemma}
\begin{proof}
    \uses{lem:barycenter-dist-bound}
    The diameter of a chain simplex is the longest of its edges
    $[\,b(F),\,b(G)\,]$ with $F \subsetneq G$ inside one maximal cell of
    dimension $d \le n$. By \Cref{lem:barycenter-dist-bound} each such edge has
    length at most $\frac{d}{d+1}\,\mathrm{diam}(G) \le \frac{n}{n+1}$ times the
    old mesh (using that $t \mapsto \frac{t}{t+1}$ is increasing). Taking the
    maximum over chain simplices bounds the new mesh by $\frac{n}{n+1}$ times the
    old.
\end{proof}

\begin{lemma}[Iterated subdivision]
    \label{lem:iterated-subdivision}
    \lean{LeanEval.Topology.iterated_subdivision}
    \leanfile{LeanEval/Topology/Brouwer/Subdivision.lean}
    \uses{lem:barycentric-shrink}
    Applying barycentric subdivision $k$ times to the trivial triangulation of
    $S_n$ (whose single maximal cell is $S_n$) yields a triangulation of mesh at
    most $\left(\frac{n}{n+1}\right)^k \operatorname{diam}(S_n)$.
\end{lemma}
\begin{proof}
    \uses{lem:barycentric-shrink}
    Induct on $k$. For $k = 0$ the mesh is $\operatorname{diam}(S_n)$. The
    inductive step multiplies the bound by $\frac{n}{n+1}$ via
    \Cref{lem:barycentric-shrink}.
\end{proof}

\begin{lemma}[Fine triangulations exist]
    \label{lem:fine-triangulation}
    \lean{LeanEval.Topology.fine_triangulation}
    \leanfile{LeanEval/Topology/Brouwer/Subdivision.lean}
    \uses{lem:iterated-subdivision}
    \leanok
    For every $\varepsilon > 0$ there is a triangulation of $S_n$ with mesh less
    than $\varepsilon$.
\end{lemma}
\begin{proof}
    \uses{lem:iterated-subdivision}
    Since $\frac{n}{n+1} < 1$, the bound
    $\left(\frac{n}{n+1}\right)^k \operatorname{diam}(S_n)$ of
    \Cref{lem:iterated-subdivision} tends to $0$ as $k \to \infty$; choose $k$
    large enough that it falls below $\varepsilon$.
\end{proof}

\section{Brouwer for the simplex}

\begin{lemma}[A non-increasing positive coordinate exists]
    \label{lem:label-pigeonhole}
    \lean{LeanEval.Topology.label_pigeonhole}
    \leanfile{LeanEval/Topology/Brouwer/Simplex.lean}
    \leanok
    Let $(x_0, \dots, x_n)$ and $(y_0, \dots, y_n)$ be nonnegative tuples with
    $\sum_i x_i = \sum_i y_i = 1$. Then there is an index $i$ with $x_i > 0$ and
    $y_i \le x_i$.
\end{lemma}
\begin{proof}
    Restrict to $J = \{\, i : x_i > 0 \,\}$, nonempty since $\sum_i x_i = 1$. On
    the complement $x_i = 0$, so $\sum_{i \in J} x_i = 1 \ge \sum_{i \in J} y_i$.
    Applying \texttt{Finset.exists\_le\_of\_sum\_le} to this inequality over the
    nonempty $J$ yields some $i \in J$ with $y_i \le x_i$, and $i \in J$ means
    $x_i > 0$.
\end{proof}

\begin{definition}[Labeling induced by a self-map]
    \label{def:induced-labeling}
    \lean{LeanEval.Topology.inducedLabeling}
    \leanfile{LeanEval/Topology/Brouwer/Simplex.lean}
    \uses{def:triangulation, lem:label-pigeonhole}
    \leanok
    Let $g : S_n \to S_n$ be continuous, written in barycentric coordinates so
    that points of $S_n$ correspond to tuples $(x_0, \dots, x_n)$ with
    $x_i \ge 0$ and $\sum_i x_i = 1$. Given a triangulation, label each vertex
    $v$ by an index $i$ with $x_i(v) > 0$ and $x_i(g(v)) \le x_i(v)$; such an
    index exists by \Cref{lem:label-pigeonhole} applied to the coordinate tuples
    of $v$ and $g(v)$.
\end{definition}

\begin{lemma}[The induced labeling is Sperner]
    \label{lem:induced-is-sperner}
    \lean{LeanEval.Topology.induced_is_sperner}
    \leanfile{LeanEval/Topology/Brouwer/Simplex.lean}
    \uses{def:induced-labeling, def:sperner-labeling}
    For any continuous $g : S_n \to S_n$ and any triangulation, the induced
    labeling of \Cref{def:induced-labeling} is a Sperner labeling.
\end{lemma}
\begin{proof}
    \uses{def:induced-labeling}
    A vertex $v$ on the face carried by the index set $I$ has $x_i(v) = 0$ for
    $i \notin I$. Its induced label is an index $i$ with $x_i(v) > 0$, hence
    $i \in I$. So the labeling respects carrier faces.
\end{proof}

\begin{lemma}[Uniform coordinate modulus]
    \label{lem:coordinate-modulus}
    \lean{LeanEval.Topology.coordinate_modulus}
    \leanfile{LeanEval/Topology/Brouwer/Simplex.lean}
    Fix a continuous $g : S_n \to S_n$. There is a function
    $\omega_g : (0, \infty) \to (0, \infty)$ with $\omega_g(\varepsilon) \to 0$
    as $\varepsilon \to 0$ such that for all $p, q \in S_n$ with
    $\lVert p - q \rVert < \varepsilon$ and all $i$ one has
    $|x_i(p) - x_i(q)| \le \omega_g(\varepsilon)$ and
    $|x_i(g(p)) - x_i(g(q))| \le \omega_g(\varepsilon)$.
\end{lemma}
\begin{proof}
    Each barycentric coordinate $x_i$ is continuous, and $S_n$ is compact, so by
    Heine--Cantor (\texttt{IsCompact.uniformContinuousOn\_of\_continuous}) the maps
    $x_i$ and $x_i \circ g$ are uniformly continuous on $S_n$. Taking $\omega_g$
    to be the (finite) maximum over the $2(n+1)$ maps of their moduli of uniform
    continuity gives a single modulus with the stated bounds and
    $\omega_g(\varepsilon) \to 0$.
\end{proof}

\begin{lemma}[Rainbow cell has a witness vertex per label]
    \label{lem:rainbow-vertex-bound}
    \lean{LeanEval.Topology.rainbow_vertex_bound}
    \leanfile{LeanEval/Topology/Brouwer/Simplex.lean}
    \uses{def:induced-labeling}
    If $C$ is a rainbow cell for the labeling induced by $g$, then for each $i$
    there is a vertex $v_i$ of $C$ labeled $i$, and by construction it satisfies
    $x_i(g(v_i)) \le x_i(v_i)$.
\end{lemma}
\begin{proof}
    \uses{def:induced-labeling}
    Rainbow means the $n+1$ vertices of $C$ carry all distinct labels
    $\{0, \dots, n\}$, so for each $i$ exactly one vertex $v_i$ is labeled $i$. By
    \Cref{def:induced-labeling} a vertex labeled $i$ satisfies
    $x_i(g(v_i)) \le x_i(v_i)$.
\end{proof}

\begin{lemma}[One-sided coordinate bound on a rainbow cell]
    \label{lem:rainbow-one-sided}
    \lean{LeanEval.Topology.rainbow_one_sided}
    \leanfile{LeanEval/Topology/Brouwer/Simplex.lean}
    \uses{lem:rainbow-vertex-bound, lem:coordinate-modulus}
    Let $\omega_g$ be the modulus of \Cref{lem:coordinate-modulus}. If a
    triangulation has mesh $< \varepsilon$ and $C$ is a rainbow cell for the
    induced labeling, then every $x \in C$ satisfies
    $x_i(g(x)) \le x_i(x) + \omega_g(\varepsilon)$ for all $i$.
\end{lemma}
\begin{proof}
    \uses{lem:rainbow-vertex-bound, lem:coordinate-modulus}
    Fix $i$ and let $v_i$ be the witness vertex of \Cref{lem:rainbow-vertex-bound},
    so $x_i(g(v_i)) \le x_i(v_i)$. Since $\mathrm{diam}(C) < \varepsilon$ we have
    $\lVert x - v_i \rVert < \varepsilon$, so \Cref{lem:coordinate-modulus} gives
    $x_i(g(x)) \le x_i(g(v_i)) + \omega_g(\varepsilon)$ and
    $x_i(v_i) \le x_i(x) + \omega_g(\varepsilon)$. Chaining the three inequalities
    yields $x_i(g(x)) \le x_i(x) + 2\,\omega_g(\varepsilon)$; replacing $\omega_g$
    by $2\,\omega_g$ (still a modulus tending to $0$) gives the stated bound.
\end{proof}

\begin{lemma}[A rainbow cell is approximately fixed]
    \label{lem:rainbow-approx-fixed}
    \lean{LeanEval.Topology.rainbow_approx_fixed}
    \leanfile{LeanEval/Topology/Brouwer/Simplex.lean}
    \uses{lem:rainbow-one-sided}
    In the setting of \Cref{lem:rainbow-one-sided}, every $x$ in a rainbow cell
    $C$ satisfies
    $|x_i(g(x)) - x_i(x)| \le (n+1)\,\omega_g(\varepsilon)$ for all $i$; hence
    $C$ contains a point within $(n+1)\,\omega_g(\varepsilon)$ of being fixed.
\end{lemma}
\begin{proof}
    \uses{lem:rainbow-one-sided}
    By \Cref{lem:rainbow-one-sided}, $x_i(g(x)) - x_i(x) \le \omega_g(\varepsilon)$
    for every $i$. Since $\sum_i x_i(g(x)) = \sum_i x_i(x) = 1$, the displacements
    sum to zero, so for each fixed $i$,
    \[
        x_i(x) - x_i(g(x)) = \sum_{j \ne i}\bigl(x_j(g(x)) - x_j(x)\bigr)
        \le n\,\omega_g(\varepsilon).
    \]
    Combining the two one-sided estimates gives
    $|x_i(g(x)) - x_i(x)| \le (n+1)\,\omega_g(\varepsilon)$.
\end{proof}

\begin{lemma}[Vanishing displacement yields a fixed point]
    \label{lem:approx-fixed-limit}
    \lean{LeanEval.Topology.approx_fixed_limit}
    \leanfile{LeanEval/Topology/Brouwer/Simplex.lean}
    \uses{def:fpp}
    Let $K \subseteq E_d$ be compact and $g : K \to K$ continuous. If there is a
    sequence $(x_k)$ in $K$ with $\lVert g(x_k) - x_k \rVert \to 0$, then $g$ has
    a fixed point in $K$.
\end{lemma}
\begin{proof}
    By compactness of $K$, the sequence $(x_k)$ has a subsequence converging to
    some $x \in K$. Along it $g(x_{k_j}) \to g(x)$ by continuity, while
    $g(x_{k_j}) - x_{k_j} \to 0$, so
    $g(x) - x = \lim_j \bigl(g(x_{k_j}) - x_{k_j}\bigr) = 0$, i.e.\ $g(x) = x$.
\end{proof}

\begin{theorem}[Brouwer for the standard simplex]
    \label{thm:simplex-brouwer}
    \lean{LeanEval.Topology.simplex_brouwer}
    \leanfile{LeanEval/Topology/Brouwer/Simplex.lean}
    \uses{def:fpp, lem:sperner, lem:fine-triangulation, lem:induced-is-sperner, lem:rainbow-approx-fixed, lem:approx-fixed-limit, lem:coordinate-modulus}
    The corner simplex $S_n$ has the fixed-point property.
\end{theorem}
\begin{proof}
    \uses{lem:sperner, lem:fine-triangulation, lem:induced-is-sperner, lem:rainbow-approx-fixed, lem:approx-fixed-limit, lem:coordinate-modulus}
    Let $g : S_n \to S_n$ be continuous and let $\omega_g$ be the combined
    coordinate modulus furnished by \Cref{lem:coordinate-modulus} (all
    uniform-continuity plumbing is confined to that lemma; what follows is pure
    sequence assembly). For each $k$ choose by \Cref{lem:fine-triangulation} a
    triangulation of mesh $< 1/k$; its induced labeling is Sperner by
    \Cref{lem:induced-is-sperner}, so by \Cref{lem:sperner} it has a rainbow
    cell, and by \Cref{lem:rainbow-approx-fixed} a point $x_k$ with
    $\lVert g(x_k) - x_k \rVert \le (n+1)\,\omega_g(1/k)$. Since
    $\omega_g(1/k) \to 0$, \Cref{lem:approx-fixed-limit} applied to the compact
    $S_n$ produces a point $x$ with $g(x) = x$.
\end{proof}

\section{Reduction to the simplex}

\begin{lemma}[Full-dimensional convex bodies are homeomorphic to balls]
    \label{lem:convex-body-homeo-ball}
    \lean{LeanEval.Topology.convex_body_homeo_ball}
    \leanfile{LeanEval/Topology/Brouwer/Reduction.lean}
    Let $B \subseteq E_m$ be convex, bounded, and with nonempty interior. Then
    its closure is homeomorphic to the closed unit ball of $E_m$; in particular a
    compact convex set in $E_m$ with nonempty interior is homeomorphic to that
    ball.
\end{lemma}
\begin{proof}
    Mathlib's
    \texttt{exists\_homeomorph\_image\_interior\_closure\_frontier\_eq\_unitBall}
    provides a self-homeomorphism of $E_m$ carrying $\operatorname{interior} B$
    to the open unit ball and $\operatorname{closure} B$ to the closed unit ball.
    Restricting this homeomorphism to $\operatorname{closure} B$ gives a
    homeomorphism onto the closed unit ball of $E_m$.
\end{proof}

\begin{lemma}[A nonempty convex set has nonempty relative interior]
    \label{lem:convex-relint-nonempty}
    \lean{LeanEval.Topology.convex_relint_nonempty}
    \leanfile{LeanEval/Topology/Brouwer/Reduction.lean}
    Every nonempty convex set $K \subseteq E_d$ has nonempty intrinsic (relative)
    interior; equivalently, $K$ has nonempty interior inside its affine span.
\end{lemma}
\begin{proof}
    This is \texttt{intrinsicInterior\_nonempty} (equivalently
    \texttt{Set.Nonempty.intrinsicInterior}): for a convex set the intrinsic
    interior is nonempty as soon as the set is, and the intrinsic interior is by
    definition the interior taken within the affine span.
\end{proof}

\begin{lemma}[Reduction to full dimension]
    \label{lem:affine-span-iso}
    \lean{LeanEval.Topology.affine_span_iso}
    \leanfile{LeanEval/Topology/Brouwer/Reduction.lean}
    \uses{lem:convex-relint-nonempty}
    Let $K \subseteq E_d$ be nonempty, compact, and convex, and let $m$ be the
    dimension of its affine span $A$. Then there is an affine isometry of $A$
    onto $E_m$ under which $K$ becomes a compact convex set with nonempty
    interior (its relative interior) in $E_m$.
\end{lemma}
\begin{proof}
    \uses{lem:convex-relint-nonempty}
    The affine span $A$ is an $m$-dimensional affine subspace of $E_d$, hence
    carries an affine isometry onto $E_m$ (\texttt{AffineSubspace.isometryEquivMap}
    composed with \texttt{AffineIsometryEquiv.ofTop}, using
    \texttt{finrank\_euclideanSpace}). By \Cref{lem:convex-relint-nonempty} the
    convex set $K$ has nonempty relative interior inside $A$, and its image under
    the isometry is a compact convex subset of $E_m$ with nonempty interior.
\end{proof}

\begin{lemma}[Compact convex bodies are homeomorphic to balls]
    \label{lem:convex-homeo-ball}
    \lean{LeanEval.Topology.convex_homeo_ball}
    \leanfile{LeanEval/Topology/Brouwer/Reduction.lean}
    \uses{lem:convex-body-homeo-ball, lem:affine-span-iso}
    Let $K \subseteq E_d$ be nonempty, compact, and convex. Then the subspace
    $K$ is homeomorphic to the closed unit ball of $E_m$, where $m$ is the
    dimension of the affine span of $K$ (with the closed ball of $E_0$
    understood as a single point).
\end{lemma}
\begin{proof}
    \uses{lem:convex-body-homeo-ball, lem:affine-span-iso}
    By \Cref{lem:affine-span-iso} an affine isometry carries $K$ to a compact
    convex set with nonempty interior in $E_m$, a homeomorphism of subspaces. By
    \Cref{lem:convex-body-homeo-ball} that image is homeomorphic to the closed
    unit ball of $E_m$; composing the two gives a homeomorphism of $K$ onto the
    ball.
\end{proof}

\begin{lemma}[The closed ball has the fixed-point property]
    \label{lem:ball-fpp}
    \lean{LeanEval.Topology.ball_fpp}
    \leanfile{LeanEval/Topology/Brouwer/Reduction.lean}
    \uses{def:fpp, thm:simplex-brouwer, lem:fpp-homeo, lem:convex-homeo-ball}
    For every $m$, the closed unit ball of $E_m$ has the fixed-point property.
\end{lemma}
\begin{proof}
    \uses{thm:simplex-brouwer, lem:fpp-homeo, lem:convex-homeo-ball}
    For $m = 0$ the ball is a single point and trivially has FPP. For $m \ge 1$,
    the corner simplex $S_m$ is a nonempty compact convex body, hence by
    \Cref{lem:convex-homeo-ball} homeomorphic to the closed unit ball of $E_m$.
    As $S_m$ has FPP by \Cref{thm:simplex-brouwer}, the closed ball has FPP by
    \Cref{lem:fpp-homeo}.
\end{proof}

\begin{lemma}[Compact convex sets have the fixed-point property]
    \label{lem:convex-fpp}
    \lean{LeanEval.Topology.convex_fpp}
    \leanfile{LeanEval/Topology/Brouwer/Reduction.lean}
    \uses{def:fpp, lem:convex-homeo-ball, lem:ball-fpp, lem:fpp-homeo}
    Every nonempty compact convex set $K \subseteq E_d$, as a subspace, has the
    fixed-point property.
\end{lemma}
\begin{proof}
    \uses{lem:convex-homeo-ball, lem:ball-fpp, lem:fpp-homeo}
    By \Cref{lem:convex-homeo-ball}, $K$ is homeomorphic to a closed unit ball,
    which has FPP by \Cref{lem:ball-fpp}; transport along the homeomorphism via
    \Cref{lem:fpp-homeo}.
\end{proof}

\begin{theorem}[Brouwer fixed-point theorem]
    \label{thm:brouwer}
    \lean{LeanEval.Topology.brouwer_fixed_point}
    \leanfile{LeanEval/Topology/Brouwer.lean}
    \uses{lem:convex-fpp, lem:fpp-ambient}
    Let $K \subseteq E_d$ be nonempty, compact, and convex, and let
    $f : E_d \to E_d$ be continuous on $K$ with $f(K) \subseteq K$. Then $f$ has
    a fixed point in $K$: there exists $x \in K$ with $f(x) = x$.
\end{theorem}
\begin{proof}
    \uses{lem:convex-fpp, lem:fpp-ambient}
    By \Cref{lem:convex-fpp} the subspace $K$ has the fixed-point property, and
    \Cref{lem:fpp-ambient} converts this into an ambient fixed point for $f$,
    given that $f$ is continuous on $K$ and maps $K$ into $K$.
\end{proof}

\end{document}
